% ============================================
% MEM-05 (ZH)
% 中文翻译 — SPEC-935-R356
% ============================================
\section*{\textbf{MEM-05} --- 黄眼}
\markboth{MEM-05 --- 黄眼}{}

\noindent \textit{再也没有闭上过的东西}

\noindent \textbf{再也没有闭上过的东西}

\vspace{0.5em}\noindent\rule{\textwidth}{0.5pt}\vspace{0.5em}

\noindent 第二天早上，母亲看着我说：

\noindent \textit{“若阿金，你的眼睛。”}

\noindent 我摸了摸脸。右眼瞪得大大的。眨不了。眼皮是硬的，仿佛被从上面缝住了。

\noindent 我跑到水罐边，俯身去看。水很脏、很浑，但勉强能看见倒影。

\noindent 右眼的虹膜是黄色的。

\noindent 不是黄疸的那种黄——那会让整只眼睛变成蛋黄的颜色。这是一种浓烈、发亮的黄，黑色的瞳孔似乎比应有的更深。仿佛我的眼睛里有一口井。

\noindent 母亲在胸前画了个十字。

\noindent \textit{“你看见了什么？”}

\noindent 我没有回答。

\noindent \textit{“若阿金，看在老天份上，你在那个坑里看见了什么？”}

\noindent \textit{“什么也没有，”我撒了谎，“我什么都没看见。”}

\noindent 她知道那是谎言。但没有再问。那天，士兵们来带走女人。他们说要把姑娘们“重新安置”到海边去。母亲一直在咳，咳个不停，可他们还是把她带走了。

\noindent 卡车开走时，她回头看着我。

\noindent 她的眼睛在说：\textit{别死。}

\noindent 我没有死。

\noindent 但我宁愿当时就死了。

\vspace{0.5em}\noindent\rule{\textwidth}{0.5pt}\vspace{0.5em}

\noindent 那天夜里，我回到坑边。

\noindent 怪物在那里。这一次，它没有在吃。它躺在尸体中间，望着天空。看见我，它坐了起来。

\noindent \textit{“我就知道你会回来。”}

\noindent 我沉默着。

\noindent \textit{“过来。”}

\noindent 我过去了。这一次，我没有坐在坑边。我下去了。我的脚陷进尸堆里。有些尸体是软的。有些是硬的。有一具伸着手臂，指向虚无。我不小心踩到它的手，手发出一声干响折断了。

\noindent 靠近了，怪物身上的气味不一样了。不只是肉。是浓稠的血、湿土，还有一丝像牲畜圈气味的甜。

\noindent \textit{“你叫什么名字？”它问。}

\noindent \textit{“若阿金。”}

\noindent \textit{“若阿金。”它像在品尝这个词一样重复了一遍。“圣人的名字。若瑟、若阿金、玛利亚。生来受苦之人的名字。”}

\noindent \textit{“你呢？”}

\noindent 它笑了。笑声是一种潮湿的声音，像水从脏管道里流过。

\noindent \textit{“我没有名字。但我有饥饿。而饥饿需要一个可以呼唤的名字。”}

\noindent \textit{“就叫眼睛吧。”}

\noindent 它歪了歪头。

\noindent \textit{“为什么？”}

\noindent \textit{“因为我最先看见的，是你的一部分。”}

\noindent 它笑了。牙齿在黑暗中发亮。

\noindent \textit{“那我就叫眼睛。对你而言。”}

\noindent 它伸出手。那只巨大的、手指太多的手，碰了碰我的脸。

\noindent \textit{“你的眼睛现在和我的一样了，”它说，“你看见我所看见的。”}

\noindent \textit{“你看见什么？”}

\noindent \textit{“我看见饥饿。”}

\noindent \textit{“在谁体内？”}

\noindent \textit{“在所有人体内。但大多数人不知道。你将会知道。你将在每一张脸上看见他藏起来的东西。你将看见每个人体内的洞。你将知道是什么填满了那个洞。”}

\noindent \textit{“什么？”}

\noindent 它像鸟一样歪过头。

\noindent \textit{“我。”}

\vspace{0.5em}\noindent\rule{\textwidth}{0.5pt}\vspace{0.5em}

\noindent 第三天夜里，怪物让我看见我将会成为什么。

\noindent 它用一根手指碰了碰我的额头——一根长长的手指，指甲漆黑——然后我\textit{看见了}。

\noindent 我看见一间屋子。一间大屋子，摆满了床。躺着的男人。喊叫的男人。一动不动的男人。我看见一个男孩——比我大，但仍然是男孩——坐在角落里。男孩动着嘴，却没有声音。他在一个本子上写着字。那笔迹是我的。

\noindent \textit{“那是你，”怪物说，“几年以后。”}

\noindent \textit{“那是哪里？”}

\noindent \textit{“一个收留像我这样的人的地方。一个被社会活埋之人的地方。你将在那里待很久。”}

\noindent \textit{“那你呢？”}

\noindent \textit{“我跟你一起去。”}

\noindent 它碰了碰我的胸口。手指穿过皮肤、骨头——不疼——碰到了我体内的某个东西。

\noindent \textit{“我已经在这里了。我不会再离开。”}

\noindent \textit{“如果我不愿意呢？”}

\noindent 它又笑了。

\noindent \textit{“你没有选择，若阿金。你坐在我的坑边的那一夜，你邀请了我。不是用言语。是用饥饿。你的饥饿呼唤了我。而我来了。”}

\noindent \textit{“我可以离开。”}

\noindent \textit{“可以。但我就在你体内。”}

\noindent \textit{“如果我死了呢？”}

\noindent 它安静了。第一次，它的沉默不一样了。不是威胁。是……悲伤。

\noindent \textit{“如果你死了，我就等。”}

\noindent \textit{“等什么？”}

\noindent \textit{“等到有人想起你。打开你写的东西。读它。”}

\noindent \textit{“然后呢？”}

\noindent 它黄色的眼睛在黑暗中发光。

\noindent \textit{“然后我再走一趟。”}

\vspace{0.5em}\noindent\rule{\textwidth}{0.5pt}\vspace{0.5em}

\noindent 当我抬起头时，一只猫头鹰停在一棵枯树上注视着我。

\noindent 它没有发出声响。它就在那里，一动不动，用没有眼皮的黑眼睛看着我。它的头缓缓转动——一百八十度——然后盯住了怪物。

\noindent 怪物没有抬头。但它笑了。

\noindent \textit{“它是来认识你的。”}

\noindent \textit{“它是谁？”}

\noindent \textit{“它是我在人间的舌头。你说的一切，它都会带走。它看见的一切，我都知道。”}

\noindent 猫头鹰张开翅膀，飞走了。飞翔无声，被夜色吸收，仿佛黑暗是一块湿布。

\noindent \textit{“它会跟着我吗？”}

\noindent \textit{“不。它会在前面。铺路。”}

\vspace{0.5em}\noindent\rule{\textwidth}{0.5pt}\vspace{0.5em}

\noindent 我72岁了。我在精神病院的一间病房里写下这些。医生说我是疯子。说怪物是我头脑的产物。说黄眼是没治好的白内障。

\noindent 我知道不是。

\noindent 它现在就在这里。在我写的时候。我感觉到它压在我的后颈上。那股甜的气息。它正在读我写的东西。它正在笑。

\noindent \textit{“告诉他，”它低语，“告诉他，日记不会在你死后结束。”}

\noindent 我很快就要死了。我知道。奥利维拉医生说我的心脏撑不住了。

\noindent 但日记会继续。

\noindent 写的人不是我。

\noindent 是它。

\noindent \textit{--- 若阿金·安东尼奥·科雷亚，患者1,260。最后记录。}

\vspace{0.5em}
\noindent\textit{\textbf{↪ 路线:} \rota{MEM-04} \rotasep \rota{MEM-12} \rotasep \rota{DOC-01}}
